\documentclass[12pt,a4paper]{article}

% ============================================================
% КОДИРОВКА — XeLaTeX
% ============================================================
\usepackage{fontspec}
\usepackage{polyglossia}
\setdefaultlanguage{russian}
\setotherlanguage{english}
\setmainfont{CMU Serif}
\setsansfont{CMU Sans Serif}
\setmonofont{CMU Typewriter Text}

% ============================================================
% МАТЕМАТИКА И ОФОРМЛЕНИЕ
% ============================================================
\usepackage{amsmath,amssymb,amsthm}
\usepackage{graphicx}
\usepackage{booktabs}
\usepackage{hyperref}
\hypersetup{colorlinks, urlcolor=blue, citecolor=blue}

\title{\textbf{Модель MEML 9.6: Ковариантная нормировка орбифолда $T^2/\mathbb{Z}_3$,\\
безразмерная свёртка bulk-профилей и аксионный реликт}\\
\large С дополнением: суперадиантность и чёрные дыры}

\author{\textbf{Юров Д. В.}\\2026}
\date{\today}

\begin{document}

\maketitle

\begin{abstract}
Представлена замкнутая модификация аксионной модели MEML~9.6 в шестимерном ($6D$) пространстве-времени $M_4 \times T^2/\mathbb{Z}_3$. Введён ковариантный КТП-аппарат, в котором гравитационный сектор и комплексное скалярное поле $\Phi$ скрытого сектора классифицированы как объемные поля (bulk fields). Размерностная редукция секторов выполнена с симметричным учётом фактора дискретной факторизации орбифолда, определяющего площадь фундаментальной области как $V_2/3$. На основе безразмерного отношения интегралов волновых функций определены фундаментальный масштаб 6D-гравитации $M_6 \approx 1.45 \times 10^{17}$~ГэВ и эффективный объемный масштаб скалярного сектора $M_* \approx 4.16 \times 10^{13}$~ГэВ. Это позволило зафиксировать константу распада $f_a = \frac{1}{\sqrt{3}}M_*^2/m_{KK} = 10^{12}$~ГэВ. Применение ангармонической поправки Лита--Тёрнера при $\theta_i \approx 0.9483$ воспроизводит плотность тёмной материи $\Omega_a h^2 = 0.120017$ (Planck 2018). Обсуждается связь с чёрными дырами: суперадиантность и аксионные звёзды с критической массой $M \sim 10^{-5} M_\odot$.
\end{abstract}

% ============================================================
\section{Введение}
% ============================================================
Предыдущие Калуца--Клейновские ревизии модели (v8.1--8.3) оказались феноменологически несостоятельными из-за ограничений криогенных детекторов XENONnT и LUX-ZEPLIN (LZ), полностью исключивших тепловых КК-бозонов с древесным сечением рассеяния $\sigma_{SI} \sim 10^{-44} \ \text{см}^2$. Переход к аксионной парадигме в версиях 9.2--9.3 вывел модель из-под удара жидкого ксенона. Версия 9.4 ввела фактор факторизации площади тора $1/3$ асимметрично, а ревизия 9.5 содержала размерностную ошибку в интеграле свёртки профилей.

Настоящая финальная ревизия 9.6 полностью устраняет данные дефекты и добавляет анализ связи аксиона с физикой чёрных дыр.

% ============================================================
\section{Стабилизация модулей и статус bulk-константы $\alpha_G$}
% ============================================================
Характерный масштаб компактификации инфракрасного (IR) сектора $m_{KK}$ определяется вакуумным ожиданием радиона:
\begin{equation}
R_{T^2} \approx \sqrt{\mathcal{I}} \cdot M_{Pl}^{-1} \exp\left(\frac{2\pi}{3\alpha_G}\right),
\label{eq:radion}
\end{equation}
где $M_{Pl} = 1.22 \times 10^{19}$~ГэВ, а $\mathcal{I} = 0.3653$ --- модулярный интеграл нулевой моды гравитона.

Для $R_{T^2} \approx 1.97 \times 10^{-31}$~м ($m_{KK} = 10^{15}$~ГэВ) bulk-константа связи $SU(5)$ составляет $\alpha_G \approx 0.2114$. В модели MEML 9.6 этот параметр честно признаётся \textbf{свободным феноменологическим постулатом}, указывающим на режим сильной связи в многомерном bulk-пространстве.

% ============================================================
\section{Ковариантная орбифолдная редукция и разделение масштабов}
% ============================================================
И многомерный метрический тензор гравитации, и комплексное скалярное поле $\Phi$ являются объемными полями. Интегрирование выполняется по фундаментальной области орбифолда $T^2/\mathbb{Z}_3$ с площадью $V_{\text{orb}} = V_2/3$.

\subsection{Безразмерная свёртка профилей и подавление амплитуды}
Отношение констант извлекается через безразмерный функционал участия:
\begin{equation}
\left(\frac{M_*}{M_6}\right)^4 = \frac{\int_{T^2/\mathbb{Z}_3} |\Psi_{\Phi}(z,\bar{z})|^4 \, d^2z}{\left( \int_{T^2/\mathbb{Z}_3} |\Psi_{\Phi}(z,\bar{z})|^2 \, d^2z \right)^2} \cdot V_{\text{orb}}.
\label{eq:convol}
\end{equation}
Фактор подавления $\epsilon \approx 2.861 \times 10^{-4}$. Масштабы:
\begin{equation}
M_{Pl}^2 = \frac{M_6^4}{3 \cdot m_{KK}^2} \implies M_6 = \left(3 M_{Pl}^2 \cdot m_{KK}^2\right)^{1/4} \approx 1.45 \times 10^{17} \ \text{ГэВ}.
\end{equation}
\begin{equation}
M_* = M_6 \cdot \epsilon \approx 4.162 \times 10^{13} \ \text{ГэВ}.
\end{equation}

Константа распада Печчеи--Квинн:
\begin{equation}
\boxed{\; f_a = \frac{1}{\sqrt{3}} \frac{M_*^2}{m_{KK}} = 10^{12} \ \text{ГэВ}. \;}
\end{equation}
Иерархия $M_6 / M_* \approx 3500$.

% ============================================================
\section{Ангармоническая релаксация и космологические струны}
% ============================================================
Масса аксиона генерируется QCD-инстантонами:
\begin{equation}
m_a = \frac{\Lambda_{\text{QCD}}^2}{f_a} = \frac{(0.075 \ \text{ГэВ})^2}{10^{12} \ \text{ГэВ}} \approx 5.63 \times 10^{-6} \ \text{эВ}.
\end{equation}

При $T_{\text{reheating}} < f_a = 10^{12}$~ГэВ космологические струны подавлены. Эволюция плотности энергии с учётом поправки Лита--Тёрнера:
\begin{equation}
\Omega_a h^2 \approx 0.120 \cdot \left(\frac{f_a}{10^{12} \ \text{ГэВ}}\right)^{7/6} \cdot \theta_i^2 \cdot f_{\text{anh}}(\theta_i),
\label{eq:omega_final}
\end{equation}
\begin{equation}
f_{\text{anh}}(\theta_i) = \left[ \ln \left( \frac{e}{1 - (\theta_i/\pi)^2} \right) \right]^{7/6}.
\end{equation}

При $\theta_i \approx 0.9483$:
\begin{equation}
\boxed{\; \Omega_a h^2 = 0.120017 \approx 0.120. \;}
\end{equation}
Совпадение с Planck 2018 в пределах $0.013\%$.

% ============================================================
\section{Честный статус свободных параметров теории}
% ============================================================
Три свободных макро-параметра:
\begin{itemize}
\item $\alpha_G = 0.2114$ --- bulk-константа связи;
\item $\epsilon = 2.861 \times 10^{-4}$ --- фактор свёртки;
\item $\theta_i = 0.9483$ --- начальный угол смещения.
\end{itemize}

Уникальная резонансная частота сигнала:
\begin{equation}
\nu_a = \frac{m_a c^2}{h} \approx 1.36 \ \text{ГГц}.
\end{equation}

% ============================================================
\section{Экспериментальная верификация в галоскопах}
% ============================================================
Аксионная природа тёмной материи снимает проблему прямого детектирования ($\sigma_{SI} = 0$). Проверка — в микроволновых полостных галоскопах.

\begin{table}[h!]
\centering
\caption{Чувствительность мировых криогенных галоскопов к сигналу MEML 9.6.}
\label{tab:haloscopes}
\begin{tabular}{lccc}
\toprule
\textbf{Эксперимент} & \textbf{Диапазон (ГГц)} & \textbf{Статус} & \textbf{Чувствительность} \\
\midrule
\textbf{ADMX (США)} & $0.6$--$2.0$ & Работает & Прямое сканирование \\
HAYSTAC (США) & $4.0$--$6.0$ & Работает & Выше диапазона \\
MADMAX (Германия) & $10$--$100$ & В разработке & Тяжёлые аксионы \\
\midrule
\textbf{MEML 9.6} & $\mathbf{1.36}$ & Теория & Проверка в ADMX \\
\bottomrule
\end{tabular}
\end{table}

% ============================================================
\section{Связь с чёрными дырами}
\label{sec:black-holes}
% ============================================================

\subsection{Суперадиантность}

Аксион MEML~9.6 может взаимодействовать с вращающимися чёрными дырами Керра через эффект \textbf{суперадиантности} (Пенроуз, 1971; Зельдович, 1972). Вращающаяся чёрная дыра усиливает волны, проходящие вблизи эргосферы, извлекая энергию из собственного вращения.

Условие усиления:
\begin{equation}
\omega < m \cdot \Omega_H,
\end{equation}
где $\Omega_H$ --- угловая скорость горизонта:
\begin{equation}
\Omega_H = \frac{a_*\,c^3}{2GM_{\text{BH}}\left(1 + \sqrt{1-a_*^2}\right)}.
\end{equation}

Вводя «гравитационную постоянную тонкой структуры»:
\begin{equation}
\alpha \equiv \frac{G M_{\text{BH}} m_a}{\hbar c},
\end{equation}
получаем, что максимум суперадиантности достигается при $\alpha \sim 0.5$. Для аксиона MEML~9.6 это даёт критическую массу:
\begin{equation}
\boxed{\; M_{\text{BH}} \sim \frac{0.5\,\hbar c}{G\,m_a}
\approx 1.2 \times 10^{-5}\,M_\odot. \;}
\end{equation}

Такие чёрные дыры должны \textbf{терять вращение} --- проверяемое предсказание для LIGO/Virgo.

\subsection{Аксионные звёзды}

Гравитационно связанное облако аксионов формирует \textbf{аксионную звезду} с критической массой:
\begin{equation}
M_{\text{crit}} \sim \frac{M_{Pl}^2}{m_a}
\approx 1.6 \times 10^{-5}\,M_\odot.
\end{equation}

\textbf{Замечательно:} эта масса совпадает с суперадиантностной. При $M > M_{\text{crit}}$ аксионная звезда коллапсирует в чёрную дыру, создавая цикл:
\begin{equation}
\text{Аксион} \to \text{Облако} \to \text{ЧД}
\to \text{Хокинг} \to \text{Аксион}.
\end{equation}

\subsection{Наблюдательные предсказания}

\begin{table}[h!]
\centering
\small
\caption{Предсказания MEML 9.6 для наблюдений чёрных дыр.}
\label{tab:bh-predictions}
\begin{tabular}{@{}p{3.5cm}p{4cm}p{4cm}@{}}
\toprule
\textbf{Наблюдение} & \textbf{Предсказание} & \textbf{Инструмент} \\
\midrule
Спин ЧД & $a_* \lesssim 0.1$ при $M \sim 10^{-5} M_\odot$ & LIGO O4/O5 \\
Тени ЧД & Искажения от облака & EHT-next \\
Стохастический фон ГВ & $\Omega_{\text{GW}} \sim 10^{-9}$ & PTA (NANOGrav) \\
\bottomrule
\end{tabular}
\end{table}

% ============================================================
\section{Выводы и открытые проблемы}
% ============================================================
\subsection*{Ключевые результаты:}
\begin{enumerate}
\item Введена безразмерная нормировка объёмного интеграла свёртки профилей мод.
\item Выведен масштаб bulk-гравитации $M_6 = 1.45 \times 10^{17}$~ГэВ.
\item Зафиксирована $f_a = 10^{12}$~ГэВ при $M_* = 4.162 \times 10^{13}$~ГэВ.
\item Плотность $\Omega_a h^2 = 0.120017$ согласуется с Planck 2018.
\item Предсказана суперадиантность для ЧД с $M \sim 10^{-5} M_\odot$.
\item Критическая масса аксионных звёзд $M_{\text{crit}} \sim 1.6 \times 10^{-5} M_\odot$.
\end{enumerate}

\subsection*{Открытые проблемы:}
\begin{enumerate}
\item Вычисление петельных поправок к $g_{a\gamma\gamma}$.
\item Анализ стабильности конических дефектов тора.
\item Согласование с полными уравнениями Эйнштейна.
\item Вывод $\alpha_G = 0.2114$ из первых принципов $SU(5)$.
\item Численное моделирование суперадиантного роста.
\end{enumerate}

\begin{thebibliography}{9}
\bibitem{Planck2018} Planck Collaboration. // Astron. Astrophys. \textbf{641}, A6 (2020).
\bibitem{PecceiQuinn1977} Peccei R.~D., Quinn H.~R. // Phys. Rev. Lett. \textbf{38}, 1440 (1977).
\bibitem{Preskill1983} Preskill J., Wise M.~B., Wilczek F. // Phys. Lett. B \textbf{120}, 127 (1983).
\bibitem{LythTurner} Lyth D.~H., Turner M.~S. // Phys. Rev. D \textbf{48}, 1502 (1993).
\bibitem{ADMX2025} ADMX Collaboration. // Phys. Rev. Lett. \textbf{134}, 021002 (2025).
\bibitem{Penrose1971} Penrose R. // Nuovo Cimento Lett. \textbf{1}, 252 (1971).
\bibitem{Zeldovich1972} Зельдович Я.~Б. // Письма в ЖЭТФ \textbf{14}, 270 (1971).
\bibitem{PressTeukolsky1972} Press W.~H., Teukolsky S.~A. // Nature \textbf{238}, 211 (1972).
\bibitem{NANOGrav2023} NANOGrav Collaboration. // Astrophys. J. Lett. \textbf{951}, L8 (2023).
\end{thebibliography}

\end{document}